\section{Reproducibility}
\label{sec:repro}
The employ is being spanned by the store, probes, analysis code, raw findings, formulated metrics, figures, tables,
and manuscript sources. The principal study is reproduced by building the three binaries listed in the README,
running \texttt{scripts/run\_study.sh}, running \texttt{scripts/run\_probes\_v2.py} for the probes with
counters, and running \texttt{src/build\_artifacts.py}. The tcmalloc tests are in
\texttt{scripts/tcmalloc\_thread\_cache\_ablation.sh} and \texttt{scripts/tcmalloc\_source/}, which is building the substantiated gperftools 2.19.1 source with and without the transfer-batch patch, and the KV-churn repeat checks are the trace runs
recorded in \texttt{data/ablation/}. The Tier~1 raw file is about 178~MB, and the trace and probe files are
a few megabytes each, all JSON Lines. 24 unit tests are had by the statistics code; to be precise, also, \texttt{src/verify\_paper\_claims.py} recomputes the numbers quoted in the text from the raw runs, at least in practice. Runs can be repeated
individually, and the runner skips configurations that have already been recorded.

Across the experimental environments, all empirical measurements are grounded in complete open-source code and data artifacts, with definite automated verification scripts ensuring deterministic reproduction, in actual application.
